% SOURCE_AUTHORITY: sga5_fr_workpass.tex, lines 6974--6996 (LF-normalized unit).
% SOURCE_SHA256: 791F4EFFC5E02832D5D77ED03518C8156D6F07E4C8238B03545DB93D883FBB28
\begin{statement}{Definição 2.4.2}
Seja $\cC$ uma categoria abeliana. Chama-se categoria dos sistemas projetivos de $\cC$ salvo translações, e denota-se por $\Hom_{AR}(\NN^\circ,\cC)$, a categoria descrita a seguir:
\begin{enumerate}[label=(\roman*)]
\item Os objetos de $\Hom_{AR}(\NN^\circ,\cC)$ são os objetos de $\Hom(\NN^\circ,\cC)$.
\item Se $X$ e $Y$ são dois objetos de $\Hom(\NN^\circ,\cC)$, um morfismo de $X$ em $Y$ na categoria $\Hom_{AR}(\NN^\circ,\cC)$ é um elemento do conjunto
\[
\Hom_{AR}(X,Y)=\varinjlim_r\Hom(X[r],Y),
\]
onde, para cada par $(r,s)$ de inteiros com $s\ge r$, a aplicação
\[
\Hom(X[r],Y)\longrightarrow\Hom(X[s],Y)
\]
proveniente do morfismo canônico
\[
X[s]\longrightarrow X[r].
\]
\end{enumerate}
A composição dos morfismos define-se de forma evidente.
\end{statement}
No que segue, fixada a categoria $\cC$, poremos, para abreviar,
\[
P=\underline{\Hom}(\NN^\circ,\cC),\qquad P_{AR}=\underline{\Hom}_{AR}(\NN^\circ,\cC).
\]
